% ================================================================
%  §6  Smooth Unconstrained Optimization
% ================================================================
\section{Smooth Unconstrained Optimization}\label{sec:smooth-methods}

The quadratic case gave us clean formulas and exact stepsizes.
Real objectives are messier: no finite termination, no closed-form line search, and convergence rates that depend on properties we may not know.
This section extends the gradient method to general smooth functions, develops the line search machinery needed to make it work, and then introduces second-order and quasi-Newton methods that converge faster --- at a price.
The unifying view: each algorithm constructs a \textit{model} of $f$ (first-order, second-order, or data-driven) and uses the model to choose the next iterate.
The model is always wrong to some extent, and managing that mismatch is the core challenge.


% ────────────────────────────────────────────────────────────────
\subsection{The Gradient Method}\label{ssec:gradient-general}

The algorithm is the same as for quadratics, but now the tomography $\varphi_{x^i, d^i}(\alpha) = f(x^i + \alpha d^i)$ is no longer a parabola --- its shape depends on $f$.

\begin{algorithm}
	\caption{Gradient Method with Steepest Descent}
	\begin{algorithmic}[1]
		\Require $f$, initial $x$, tolerance $\varepsilon > 0$
		\Ensure Approximate stationary point of $f$
		\Procedure{SD}{$f, x, \varepsilon$}
			\While{$\norm{\nabla f(x)} > \varepsilon$}
				\State $d \gets -\nabla f(x)$
				\State $\alpha \gets \mathrm{stepsize}(f, x, d)$
				\State $x \gets x + \alpha d$
			\EndWhile
		\EndProcedure
	\end{algorithmic}
\end{algorithm}

At each iterate $x^i$, the first-order model $L^i(x) = f(x^i) + \inner{\nabla f(x^i), x - x^i}$ is a good approximation only in a small ball around $x^i$.
The update $x^{i+1} = x^i + \alpha^i d^i$ is meaningful only when $f(x^{i+1}) < f(x^i)$.

The stepsize $\alpha^i$ is the critical design choice.
Too large and the method overshoots: $f(x^{i+1}) > f(x^i)$.
Too small and the method stagnates, decreasing $f$ only infinitesimally per step.
Scylla and Charybdis --- and finding a path between them is the subject of the next subsections.

Using the unnormalized gradient $d^i = -\nabla f(x^i)$ (rather than $d^i = -\nabla f(x^i)/\norm{\nabla f(x^i)}$) ties the step length to the gradient magnitude: far from the optimum the steps are large, near it they shrink automatically.


% ────────────────────────────────────────────────────────────────
\subsection{Fixed Stepsize}\label{ssec:fixed-stepsize}

A fixed stepsize $\alpha^i = \bar{\alpha}$ is simple and cheap, but it needs a guarantee that $f$ decreases at every step.
Such a guarantee requires controlling how fast $\nabla f$ can change --- i.e., $L$-smoothness (\Cref{def:L-smooth}).

The tomography along the gradient direction gives $\varphi'_{x,-\nabla f(x)}(0) = -\norm{\nabla f(x)}^2 < 0$: the initial descent is always present.
$L$-smoothness ensures this descent doesn't vanish too quickly.
If $f$ is $L$-smooth, $\varphi$ inherits $[L\norm{d}^2]$-smoothness, giving the upper bound
\begin{equation}\label{eq:tomography-upper-bound}
	\varphi(\alpha) \leq \varphi(0) + \norm{\nabla f(x)}^2 \left[\frac{L\alpha^2}{2} - \alpha\right].
\end{equation}
The worst-case quadratic bound says: the function can curve \emph{up} at rate $L$, but we know the initial slope is $-\norm{\nabla f(x)}^2$.

Choosing $\bar{\alpha} = 1/L$ yields the descent bound
\begin{equation}\label{eq:fixed-step-descent}
	f(x^{i+1}) - f(x^i) \leq -\frac{\norm{\nabla f(x^i)}^2}{2L}.
\end{equation}
The decrease is \emph{additive}, not multiplicative --- sublinear convergence:
\begin{equation}\label{eq:sublinear-smooth}
	f(x^k) - f^* \leq \frac{2L\norm{x^0 - x_*}^2}{k - 1}
	\qquad \Longrightarrow \qquad
	k \in \mathcal{O}\!\left(\frac{LD^2}{\varepsilon}\right).
\end{equation}
This matches the quadratic case when $\lambda_n = 0$ (\Cref{thm:sublinear-quad}).

The $\mathcal{O}(1/\varepsilon)$ bound is not optimal for the class of $L$-smooth convex functions: $\mathcal{O}(1/\sqrt{\varepsilon})$ is achievable (see \Cref{ssec:accelerated}).
But a lower bound shows that some problems are genuinely hard: there exist $L$-smooth functions for which any first-order algorithm requires $\Omega(1/\sqrt{\varepsilon})$ iterations.

\paragraph{Linear convergence with strong convexity}

If $f$ is also $\tau$-strongly convex, the mean value theorem gives $\nabla f(x) - \nabla f(x_*) = \nabla^2 f(w)(x - x_*)$ for some $w$.
Setting $z = x - \alpha \nabla f(x)$:
\begin{equation}\label{eq:strong-convex-contraction}
	z - x_* = (I - \alpha \nabla^2 f(w))(x - x_*),
	\qquad
	\norm{z - x_*} \leq \norm{I - \alpha \nabla^2 f(w)} \cdot \norm{x - x_*}.
\end{equation}
The contraction factor is $r = \norm{I - \alpha \nabla^2 f(w)} = \max\{|1 - \alpha\lambda_1|,\, |1 - \alpha\lambda_n|\}$.
Since $\tau \leq \lambda_n \leq \lambda_1 \leq L$, the optimal choice $\alpha = 2/(L + \tau)$ gives
\begin{equation}\label{eq:optimal-fixed-rate}
	r = \frac{L - \tau}{L + \tau} = \frac{\kappa - 1}{\kappa + 1} < 1,
	\qquad \kappa = L / \tau.
\end{equation}
Linear convergence, with rate governed by the condition number --- the same story as for quadratics, but now for any $\tau$-strongly convex, $L$-smooth objective.


% ────────────────────────────────────────────────────────────────
\subsection{Inexact Line Search}\label{ssec:line-search}

Exact line search ($\alpha^* \in \argmin_\alpha \varphi(\alpha)$) is rarely practical for general $f$.
We need cheaper alternatives that still guarantee convergence.

\paragraph{Armijo condition (avoiding Scylla: overshooting)}

\begin{equation}\label{eq:armijo}
	\varphi(\alpha) \leq \varphi(0) + m_1 \alpha\, \varphi'(0),
	\qquad 0 < m_1 \ll 1.
\end{equation}
This ensures \textit{sufficient decrease}: the function drops by at least a fraction $m_1$ of what the linear model predicts.
Small $\alpha$ always satisfies Armijo, but tiny steps cause stagnation.

\paragraph{Wolfe condition (avoiding Charybdis: stagnation)}

\begin{equation}\label{eq:wolfe}
	\varphi'(\alpha) \geq m_3\, \varphi'(0),
	\qquad m_1 < m_3 < 1.
\end{equation}
The derivative at the accepted point should be less negative than at the start --- we should have made ``enough'' progress along the ray.
The \textit{strong Wolfe condition} $|\varphi'(\alpha)| \leq m_3\, |\varphi'(0)|$ additionally prevents $\varphi'(\alpha) \gg 0$, filtering out some local maxima.

\begin{thm}[Existence of Armijo--Wolfe stepsizes]\label{thm:aw-exist}
	If $\varphi \in C^1$ and $\varphi(\alpha)$ is bounded below for $\alpha \geq 0$, then there exists $\alpha$ satisfying both Armijo and Wolfe conditions \cite{nocedal2006}.
\end{thm}

Armijo plus Wolfe captures all local minima of $\varphi$ (unless $m_1 \approx 1$).
Armijo plus strong Wolfe refines this by filtering maxima.

\paragraph{Backtracking line search}

The simplest practical approach: check only Armijo, and shrink $\alpha$ until it passes.

\begin{algorithm}
	\caption{Backtracking Line Search (BLS)}
	\begin{algorithmic}[1]
		\Require $\varphi, \alpha, m_1, \tau$ with $\tau < 1$
		\Procedure{BLS}{$\varphi, \alpha, m_1, \tau$}
			\While{$\varphi(\alpha) > \varphi(0) + m_1 \alpha\, \varphi'(0)$}
				\State $\alpha \gets \tau\, \alpha$
			\EndWhile
		\EndProcedure
	\end{algorithmic}
\end{algorithm}

Starting from $\alpha = 1$, the procedure generates $\alpha \geq \tau^h$ where $h$ is the number of shrinkage steps needed.
If stepsizes stay bounded away from zero throughout the optimization, convergence is ensured.

\paragraph{Convergence with Armijo--Wolfe}

From $L$-smoothness and the Wolfe condition, a lower bound on the accepted stepsize follows: $\alpha' > (1 - m_3)/L$.
If $f$ is also $\tau$-convex, convergence is linear with $r \approx 1 - \tau/L$, where the ``$\approx$'' depends on $m_1$ and $m_3$.


% ────────────────────────────────────────────────────────────────
\subsection{Efficiency for General Functions}\label{ssec:efficiency-general}

\begin{thm}[Local convergence rate]\label{thm:grad-general}
	Let $f \in C^2$ with $x_*$ a local minimum satisfying $\nabla^2 f(x_*) \succ 0$.
	The exact-line-search gradient method converges linearly with
	\begin{equation}\label{eq:local-convergence-rate}
		r = \left(\frac{\lambda_1 - \lambda_n}{\lambda_1 + \lambda_n}\right)^{\!2},
	\end{equation}
	where $\lambda_1, \lambda_n$ are the extreme eigenvalues of $\nabla^2 f(x_*)$.
\end{thm}

These properties need only hold in $\mathcal{B}(x_*, \delta)$.
Inexact line search may slow convergence ($r \approx 1 - \lambda_n/\lambda_1$), but reduces function evaluations per iteration --- a practical trade-off.


% ────────────────────────────────────────────────────────────────
\subsection{Twisted Gradient Methods and General Descent}\label{ssec:twisted-descent}

The direction $d^i = -\nabla f(x^i)$ is not the only descent direction.
Any $d^i$ with $\inner{d^i, \nabla f(x^i)} < 0$ works --- the entire open half-space opposite to $\nabla f$ consists of descent directions.

\begin{defn}[Descent direction]\label{def:descent-dir}
	A direction $d^i$ is a \textit{descent direction} if $\cos(\theta^i) = \inner{d^i, \nabla f(x^i)} / (\norm{d^i} \norm{\nabla f(x^i)}) < 0$.
\end{defn}

A ``twisted'' gradient method uses $d^i = H^i(-\nabla f(x^i))$ where $H^i$ is a positive definite matrix.
When $H^i \succ 0$, the direction is guaranteed to be descent.
Newton's method sets $H^i = [\nabla^2 f(x^i)]^{-1}$; quasi-Newton methods approximate this.

\begin{thm}[Zoutendijk's theorem]\label{thm:zoutendijk}
	If $f \in C^1$ is $L$-smooth with $f^* > -\infty$, and stepsizes satisfy Armijo--Wolfe, then
	\begin{equation}\label{eq:zoutendijk}
		\sum_{i=1}^{\infty} \cos^2(\theta^i)\, \norm{\nabla f(x^i)}^2 < \infty.
	\end{equation}
	If $\cos(\theta^i) \geq \bar{\theta} > 0$ for all $i$ (the angles stay bounded away from $90^\circ$), then $\norm{\nabla f(x^i)} \to 0$.
\end{thm}

Zoutendijk's theorem is the fundamental convergence guarantee for descent methods.
It says: as long as you don't search in directions nearly orthogonal to the gradient, the gradient norm must eventually vanish.


% ────────────────────────────────────────────────────────────────
\subsection{Newton's Method}\label{ssec:newton-multivariate}

Newton's method replaces the linear model with a quadratic one, using the Hessian to capture local curvature.
The payoff is quadratic convergence --- but only locally.

\begin{defn}[Newton direction]\label{def:newton-dir}
	If $\nabla^2 f(x^i) \succ 0$, the \textit{Newton direction} is
	$d^i = -[\nabla^2 f(x^i)]^{-1} \nabla f(x^i)$.
	It minimizes the second-order model $Q^i(z) = f(x^i) + \inner{\nabla f(x^i), z - x^i} + \frac{1}{2}(z - x^i)^T \nabla^2 f(x^i)(z - x^i)$.
\end{defn}

When $\nabla^2 f(x^i) \succ 0$, the Newton direction is descent: $\inner{\nabla f(x^i), d^i} = -\nabla f(x^i)^T [\nabla^2 f(x^i)]^{-1} \nabla f(x^i) < 0$.

\paragraph{Geometric interpretation}

For a quadratic $f(x) = \frac{1}{2} x^T Q x + q^T x$ with $Q \succ 0$, Newton terminates in one step.
The geometric reason: let $Q = RR^T$ and change variables to $z = Rx$.
In $z$-space the Hessian is the identity, and steepest descent converges in one step.
Newton in $x$-space \emph{is} steepest descent in the space where the Hessian is the identity.

\paragraph{Convergence}

\begin{thm}[Quadratic convergence]\label{thm:newton-quad}
	If $f \in C^3$, $\nabla f(x_*) = 0$, $\nabla^2 f(x_*) \succ 0$, then $\exists\, \delta > 0$ such that starting in $\mathcal{B}(x_*, \delta)$ gives quadratic convergence.
\end{thm}

\begin{thm}[Global convergence with line search]\label{thm:newton-global}
	If $f \in C^2$ is $L$-smooth and $\tau$-convex, then $\cos(\theta^i) \geq \tau/L$, and Zoutendijk's theorem guarantees global convergence.
	Near $x_*$, the Armijo condition with $m_1 \leq 1/2$ accepts the full Newton step ($\alpha = 1$) eventually, recovering quadratic convergence.
\end{thm}

\paragraph{The nonconvex case}

When $\nabla^2 f(x^i) \not\succ 0$, the Newton direction may not be descent.
The fix: replace $\nabla^2 f(x^i)$ with a positive definite approximation $H^i$ satisfying $\tau I \preceq H^i \preceq LI$.
A simple choice: $H^i = \nabla^2 f(x^i) + \epsilon^i I$ with $\epsilon^i = \max\{0,\, \delta - \lambda_{\min}(\nabla^2 f(x^i))\}$.
If $\{x^i\} \to x_*$ with $\nabla^2 f(x_*) \succ 0$, then $\epsilon^i = 0$ eventually, recovering quadratic convergence.
The cost is $\mathcal{O}(n^3)$ per iteration (Cholesky factorization), prohibitive for very large $n$.


% ────────────────────────────────────────────────────────────────
\subsection{Trust Region Methods}\label{ssec:trust-region}

Trust region methods reverse the line search logic: first choose a radius $r$ (the ``trust region''), then find the best direction within it.

\begin{defn}[Trust region subproblem]\label{def:trust-region}
	Given a trust region $\mathcal{B}_2(x^i, r)$:
	$x^{i+1} \in \argmin\{Q^i(z) : \norm{z - x^i} \leq r\}$.
\end{defn}

The KKT conditions for this subproblem give $[\nabla^2 f(x^i) + \lambda^i I]\, d^i = -\nabla f(x^i)$ with $\lambda^i \geq 0$.
Three cases: if $\lambda^i > 0$, the constraint is active (similar to Hessian modification with $\epsilon^i = \lambda^i$); if $\norm{d^i} < r$, then $\lambda^i = 0$ and we get the plain Newton step; as $\{x^i\} \to x_*$, eventually $\lambda^i = 0$ and quadratic convergence is recovered.

Trust region methods handle nonconvex Hessians naturally: negative-curvature directions are exploited within the trust region rather than suppressed.


% ────────────────────────────────────────────────────────────────
\subsection{Quasi-Newton Methods}\label{ssec:quasi-newton}

Computing and inverting the Hessian costs $\mathcal{O}(n^3)$.
Quasi-Newton methods approximate the Hessian (or its inverse) using only gradient differences, achieving superlinear convergence at $\mathcal{O}(n^2)$ per iteration.

\begin{defn}[Secant equation]\label{def:secant}
	Define $s^i = x^{i+1} - x^i$ and $y^i = \nabla f(x^{i+1}) - \nabla f(x^i)$.
	The \textit{secant equation} requires $H^{i+1} s^i = y^i$.
	For compatibility with $H^{i+1} \succ 0$, the \textit{curvature condition} $\inner{s^i, y^i} > 0$ must hold --- enforceable via Armijo--Wolfe line search.
\end{defn}

Superlinear convergence holds whenever the approximation resembles the true Hessian along the search direction: $\lim_{i \to \infty} \norm{(H^i - \nabla^2 f(x^i))\, d^i} / \norm{d^i} = 0$.

\paragraph{DFP formula}

The Davidon--Fletcher--Powell update finds the closest positive-definite matrix satisfying the secant equation (in a weighted Frobenius norm sense):
\begin{equation}\label{eq:dfp-update}
	H^{i+1} = (I - \rho^i y^i (s^i)^T)\, H^i\, (I - \rho^i s^i (y^i)^T) + \rho^i y^i (y^i)^T,
	\quad \rho^i = 1/\inner{s^i, y^i}.
\end{equation}

\paragraph{BFGS formula}

The Broyden--Fletcher--Goldfarb--Shanno update works on both $H^i$ and its inverse $B^i = (H^i)^{-1}$:
\begin{align*}
	H^{i+1} &= H^i + \rho^i y^i (y^i)^T - \frac{H^i s^i (s^i)^T H^i}{(s^i)^T H^i s^i}, \\
	B^{i+1} &= (I - \rho^i s^i (y^i)^T)\, B^i\, (I - \rho^i y^i (s^i)^T) + \rho^i s^i (s^i)^T.
\end{align*}
Both are rank-two updates at $\mathcal{O}(n^2)$ cost.
In practice, BFGS generally outperforms DFP.
DFP and BFGS are the endpoints of the \textit{Broyden family}: $H^{i+1} = \beta\, H^{i+1}_{\mathrm{DFP}} + (1 - \beta)\, H^{i+1}_{\mathrm{BFGS}}$ for $\beta \in [0, 1]$.

\paragraph{L-BFGS}

For very large $n$, even $\mathcal{O}(n^2)$ is too much.
L-BFGS stores only the last $k \ll n$ correction pairs $(s^i, y^i)$ and reconstructs the inverse Hessian on the fly.
Memory and time per iteration: $\mathcal{O}(kn)$.
Large $k$ approximates Newton; small $k$ degrades toward steepest descent.
A common base matrix: $B^{i-k} = \gamma^i I$ with $\gamma^i = \inner{s^{i-1}, y^{i-1}} / \norm{y^{i-1}}^2$.


% ────────────────────────────────────────────────────────────────
\subsection{Deflected Gradient Methods}\label{ssec:deflected}

Twisting ($d^i = H^i(-\nabla f(x^i))$) costs at least $\mathcal{O}(n^2)$.
\textit{Deflection} is cheaper --- $\mathcal{O}(n)$ by design:
\begin{equation}\label{eq:deflected-gradient}
	d^i = -\nabla f(x^i) + \beta^i d^{i-1}.
\end{equation}
If $\beta^0 = 0$ (starting with steepest descent), the direction aggregates all past gradients.
Ensuring descent ($\varphi'_{x^i, d^i}(0) < 0$) is nontrivial, unlike twisting where $H^i \succ 0$ suffices.


% ────────────────────────────────────────────────────────────────
\subsection{Nonlinear Conjugate Gradient}\label{ssec:nonlinear-cg}

The conjugate gradient idea from \Cref{ssec:conjugate-gradient} extends to general functions via deflection.

\begin{algorithm}
	\caption{Nonlinear Conjugate Gradient (NCG)}
	\begin{algorithmic}[1]
		\Require $f, x, \varepsilon$
		\Procedure{NCG}{$f, x, \varepsilon$}
			\State $\nabla f^- \gets 0$
			\While{$\norm{\nabla f(x)} > \varepsilon$}
				\If{$\nabla f^- = 0$}
					\State $d \gets -\nabla f(x)$ \Comment{First step: steepest descent}
				\Else
					\State $\beta \gets$ chosen formula
					\State $d \gets -\nabla f(x) + \beta\, d^-$
				\EndIf
				\State $\alpha \gets \mathrm{LS}(f, x, d)$;\; $x \gets x + \alpha d$
				\State $d^- \gets d$;\; $\nabla f^- \gets \nabla f(x)$
			\EndWhile
		\EndProcedure
	\end{algorithmic}
\end{algorithm}

Four standard formulas for $\beta$:
\begin{itemize}
	\item Fletcher--Reeves: $\beta_{\mathrm{FR}}^i = \norm{\nabla f(x^i)}^2 / \norm{\nabla f(x^{i-1})}^2$
	\item Polak--Ribi\`ere: $\beta_{\mathrm{PR}}^i = \inner{\nabla f(x^i) - \nabla f(x^{i-1}),\, \nabla f(x^i)} / \norm{\nabla f(x^{i-1})}^2$
	\item Hestenes--Stiefel: $\beta_{\mathrm{HS}}^i = \inner{\nabla f(x^i) - \nabla f(x^{i-1}),\, \nabla f(x^i)} / \inner{\nabla f(x^i) - \nabla f(x^{i-1}),\, d^{i-1}}$
	\item Dai--Yuan: $\beta_{\mathrm{DY}}^i = \norm{\nabla f(x^i)}^2 / \inner{\nabla f(x^i) - \nabla f(x^{i-1}),\, d^{i-1}}$
\end{itemize}
For quadratic $f$ with exact line search, all four reduce to the standard conjugate gradient and terminate in $n$ steps.
For general $f$, they differ.

Each formula requires specific line search conditions.
FR needs $m_1 < m_3 < 1/2$ (strong Wolfe).
PR may produce non-descent directions; the fix is $\beta_{\mathrm{PR}+}^i = \max\{\beta_{\mathrm{PR}}^i, 0\}$.
These ``$\max\{0, \cdot\}$'' modifications reset the direction to steepest descent when deflection goes wrong --- a form of \textit{restart}.
Without restarts, FR can enter a failure mode: if $\norm{\nabla f(x^i)} \ll \norm{d^i}$, then $\cos(\theta^i) \approx 0$ and the steps barely move.

In practice, all variants perform similarly, with results that vary significantly across problems.
On quadratics, $n$ conjugate gradient steps approximate one Newton step: $\norm{x^{i+n} - x_*} \leq r\, \norm{x^i - x_*}^2$.
For large $n$ this can be expensive; the method is powerful but not easy to tune.


% ────────────────────────────────────────────────────────────────
\subsection{Heavy Ball Method}\label{ssec:heavy-ball}

The heavy ball method adds \textit{momentum}: a fraction of the previous step.
\begin{equation}\label{eq:heavy-ball}
	x^{i+1} = x^i - \alpha^i \nabla f(x^i) + \beta^i (x^i - x^{i-1}).
\end{equation}
The term $\beta^i(x^i - x^{i-1})$ maintains the previous direction, reducing zig-zagging; $-\nabla f(x^i)$ steers toward $x_*$.
This is \emph{not} an $f$-descent method ($f(x^{i+1})$ may increase), but with proper $\alpha^i, \beta^i$ it achieves $d$-descent: $\norm{x^{i+1} - x_*} \leq r\, \norm{x^i - x_*}$.

The dynamics follow a two-term recurrence.
Using the mean value theorem, the state update can be written as
\begin{equation}\label{eq:heavy-ball-state}
	\begin{bmatrix} x^{i+1} - x_* \\ x^i - x_* \end{bmatrix}
	= C^i
	\begin{bmatrix} x^i - x_* \\ x^{i-1} - x_* \end{bmatrix},
	\qquad
	C^i = \begin{bmatrix}
		(1 + \beta^i)I - \alpha^i \nabla^2 f(w^i) & -\beta^i I \\
		I & 0
	\end{bmatrix}.
\end{equation}
The spectral radius $\rho(C^i)$ governs convergence.
Minimizing over $\alpha, \beta$ for the class $\tau I \preceq \nabla^2 f \preceq LI$ gives the optimal parameters:
\begin{equation}\label{eq:heavy-ball-optimal}
	\alpha = \frac{4}{(\sqrt{L} + \sqrt{\tau})^2},
	\qquad
	\beta = \left(\frac{\sqrt{L} - \sqrt{\tau}}{\sqrt{L} + \sqrt{\tau}}\right)^{\!2},
	\qquad
	r = \frac{\sqrt{\kappa} - 1}{\sqrt{\kappa} + 1}.
\end{equation}

This is the same $\sqrt{\kappa}$-dependence as conjugate gradient --- a dramatic improvement over steepest descent's $\kappa$-dependence.

\begin{defn}[Gelfand's formula and quasi-linear convergence]\label{def:gelfand}
	For any matrix $C$, $\rho(C) = \lim_{i \to \infty} \norm{C^i}^{1/i}$.
	This means: for any $\varepsilon > 0$, $\exists\, h$ such that $\norm{C^i} \leq (\rho(C) + \varepsilon)^i$ for all $i \geq h$.
	Convergence is linear if $\rho(C) + \varepsilon < 1$, but the method may not converge in the early iterates --- it needs a ``warm-up'' phase.
	This is called \textit{quasi-linear convergence}.
\end{defn}

For non-convex $f$, the method converges if $\beta \in [0, 1)$ and $\alpha \in (0, 2(1 - \beta)/L)$.
Pushing $\beta \to 1$ forces $\alpha \to 0$: there is a tension between momentum and step size.
If $\tau = 0$, the error decreases as $\mathcal{O}(1/k)$ --- same as vanilla gradient descent.


% ────────────────────────────────────────────────────────────────
\subsection{Accelerated Gradient Method}\label{ssec:accelerated}

Nesterov's accelerated gradient achieves \textit{optimal} worst-case rates among first-order methods.

\begin{algorithm}
	\caption{Accelerated Gradient (ACCG)}
	\begin{algorithmic}[1]
		\Procedure{ACCG}{$f, x, \varepsilon$}
			\State $x_- \gets x$;\; $\gamma \gets 1$
			\Repeat
				\State $\gamma_- \gets \gamma$;\;
					$\gamma \gets (\sqrt{4\gamma^2 + \gamma^4} - \gamma^2) / 2$
				\State $\beta \gets \gamma(\gamma_-^{-1} - 1)$
				\State $y \gets x + \beta(x - x_-)$
				\State $g \gets \nabla f(y)$
				\State $x_- \gets x$;\; $x \gets y - L^{-1} g$
			\Until{$\norm{g} \leq \varepsilon$}
		\EndProcedure
	\end{algorithmic}
\end{algorithm}

For $L$-smooth convex functions, the method achieves $\mathcal{O}(LD^2/k^2)$ error --- equivalently, $\mathcal{O}(1/\sqrt{\varepsilon})$ iterations.
This matches the lower bound: no first-order method can do better in the worst case.
If $f$ is also $\tau$-convex, convergence is linear with $r = (\sqrt{\kappa} - 1)/(\sqrt{\kappa} + 1)$.

The method evaluates $\nabla f$ at an \textit{extrapolated} point $y$ rather than $x$: it ``looks ahead'' before stepping.
It is designed to optimize worst-case behavior; in practice it can be slower than simpler methods on well-conditioned problems, but it achieves exactly what it is built for.